\section{Research Objectives} \label{section:research_objectives}

This project aims to evaluate and enhance the use of machine-learned interatomic potentials (MLPs) for predicting the relative stability of material interfaces. By replacing expensive density functional theory (DFT) calculations in the early stages of interface screening, MLPs offer a route to scalable evaluation across diverse chemical systems and stacking configurations.

Interfaces govern the functional behaviour of materials in microelectronics, energy storage, and optoelectronics, mediating phenomena such as band alignment, carrier confinement, and local defect formation. However, the configurational landscape of interface structures is combinatorially large. Even for lattice-matched systems, variations in Miller planes, interfacial shifts, and terminations can lead to significant changes in formation energy and electronic response. Traditional DFT methods, while accurate, are computationally prohibitive when applied to hundreds or thousands of such configurations.

This work builds on the integration of MLPs, particularly the MACE model, into high-throughput interface generation and evaluation pipelines. It focuses not only on energy prediction, but also on favourability ranking, robustness to system-specific errors, and hybrid workflows where MLPs are used as a pre-screening stage ahead of final DFT validation.

\paragraph{The research addresses the following core objectives:}

\begin{enumerate}
    \item \textbf{Benchmark MLP-predicted rankings against DFT across diverse interfaces.} Using MACE-evaluated relaxations and energies, this study quantifies how well MLPs preserve the DFT ordering of interface favourability across multiple group-IV homostructures and heterostructures.
    \item \textbf{Diagnose system-specific discrepancies and explore delta-learning corrections.} Particular attention is given to cases, such as pure-Si systems, where MLP rank fidelity degrades. Chapter~\ref{chapter:next_steps} proposes training delta models to correct these systematic energy errors.
    \item \textbf{Develop structure-based predictors of interface energy.} Surface and slab-level features (e.g. orientation, element type, symmetry, strain) will be used to train models capable of predicting favourability prior to full interface assembly, reducing reliance on supercell enumeration.
    \item \textbf{Integrate MLPs and predictors into automated workflows.} Tools such as ARTEMIS and RAFFLE will be extended to incorporate MLP-based scoring or filtering, enabling efficient identification of favourable configurations before expensive DFT refinement.
\item \end{enumerate}

In the long term, this approach supports scalable interface screening, facilitates the design of functional heterostructures, and lays groundwork for autonomous modelling pipelines. It targets regimes where standard coherence rules or bulk-derived heuristics fail, providing a data-driven alternative grounded in atomic-level energetics.
% again not the phd goal, just the goals of this project, should probably add why i aim to target those regimes.